\section{Set-of-Mark: numeración de toda la imagen para asistir la localización}

\subsection{Principio del método}

Set-of-Mark (SoM) es el trabajo representativo del Visual Prompting; el método es extremadamente simple pero muy efectivo:

\begin{enumerate}
    \item Se realiza segmentación de instancias sobre toda la imagen original (usando la función de segmentación automática de SAM)
    \item Se coloca un número en el punto central de cada región segmentada
    \item La nueva imagen con numeración, junto con la pregunta en texto, se introduce al VLM
    \item El VLM responde la pregunta de localización eligiendo un número
\end{enumerate}

\begin{importantbox}{Supuesto de diseño de SoM}
La efectividad de SoM se basa en un supuesto clave: durante el preentrenamiento, el VLM ya estuvo expuesto a una gran cantidad de labelled figures, charts y textbooks anotados. Estos datos contienen de manera natural marcas como números y flechas, por lo que el VLM tiene una fuerte capacidad de comprensión de este tipo de marcas visuales. SoM aprovecha esta capacidad: segmenta la imagen en regiones semánticas y asigna a cada región un ID legible, de modo que el VLM pueda «describir la imagen».
\end{importantbox}

\subsection{Análisis de ventajas}

Las dos ventajas centrales de SoM:

\begin{enumerate}
    \item \textbf{Reduce la dificultad de la inferencia}: convierte el problema de regresión de generar coordenadas directamente en un problema de clasificación de elegir un número
    \item \textbf{Ayuda a enfocar la atención}: las marcas numéricas tras la segmentación ayudan al modelo a enfocarse en regiones específicas, estabilizando la atención durante el proceso de inferencia
\end{enumerate}

\subsection{Resumen de la sección}

SoM, mediante la estrategia simple de segmentar toda la imagen y anotarla con números, convierte el problema de localización de objetos en una pregunta de opción múltiple, lo que reduce de manera notable la dificultad de inferencia del VLM, y constituye un ejemplo clásico de Visual Prompting.

